In 1910, the scientist Alexandre Yersin (discoverer of the bubonic plague
bacillus) established an astronomical observatory on the roof of his house in
Nha Trang, Vietnam. There, he used a Claude--Driencourt prism astrolabe,
designed to record the instant at which a star reaches a fixed zenith distance
($z = 30.0^\circ$). The following data table is inspired by Yersin's
observations under atmospheric conditions of $T = 27.0\,^\circ\mathrm{C}$ and
$P = 1011.0\,\mathrm{hPa}$.

\begin{center}
\begin{tabular}{ccccc}
\toprule
Star & Right Ascension ($\alpha$) & Hour Angle ($H$) & Recorded Time ($T_R$) & Azimuth ($A$) \\
\midrule
A & $17^{\mathrm{h}}35^{\mathrm{m}}35.55^{\mathrm{s}}$ & $+2^{\mathrm{h}}02^{\mathrm{m}}08.95^{\mathrm{s}}$ & $19^{\mathrm{h}}36^{\mathrm{m}}53.69^{\mathrm{s}}$ & $285.0^\circ$ \\
B & $19^{\mathrm{h}}59^{\mathrm{m}}52.48^{\mathrm{s}}$ & $+1^{\mathrm{h}}47^{\mathrm{m}}50.38^{\mathrm{s}}$ & $21^{\mathrm{h}}46^{\mathrm{m}}52.11^{\mathrm{s}}$ & $245.0^\circ$ \\
C & $01^{\mathrm{h}}39^{\mathrm{m}}52.31^{\mathrm{s}}$ & $-2^{\mathrm{h}}00^{\mathrm{m}}30.38^{\mathrm{s}}$ & $23^{\mathrm{h}}38^{\mathrm{m}}31.61^{\mathrm{s}}$ & $70.0^\circ$ \\
D & $01^{\mathrm{h}}41^{\mathrm{m}}38.84^{\mathrm{s}}$ & $-1^{\mathrm{h}}37^{\mathrm{m}}13.18^{\mathrm{s}}$ & $00^{\mathrm{h}}03^{\mathrm{m}}35.22^{\mathrm{s}}$ & $125.0^\circ$ \\
\bottomrule
\end{tabular}
\end{center}

\emph{Note: Note that $A$ is measured Eastward from North and $H$ is measured
Westward from South. For maintaining consistency across the problem, report
your answers to two decimal places unless otherwise specified.}

\medskip\noindent\textbf{Part 1 (15 points): The Time Problem}

\begin{description}
\item[(a)] (6 points) Calculate the difference (error) $\Delta T$ between the
  local sidereal time and the recorded time (clock time) for each star. Then,
  determine their mean $(\Delta T_{\mathrm{mean}})$ and standard deviation
  $(\sigma_T)$.
\item[(b)] (3 points) Calculate the mean of the time differences separately
  for the West stars $(\Delta T_{\mathrm{W}})$ and East stars
  $(\Delta T_{\mathrm{E}})$. Determine the discrepancy
  $\Delta T_{\mathrm{W}} - \Delta T_{\mathrm{E}}$.
\item[(c)] (6 points) The astrolabe has a systematic zenith-distance error
  $(\mathrm{d}z)$. Evaluate the uncertainty on the time difference,
  $\mathrm{d}(\Delta T)$, caused by the astrolabe. Show that by averaging the
  time differences of the East and West stars the uncertainty caused by the
  astrolabe can be reduced.
\end{description}

\medskip\noindent\textbf{Part 2 (3 points): Atmospheric Refraction}

Atmospheric refraction in radians is approximately given by
\[
R = \frac{7.888 \times 10^{-5} P}{(273.0 + T)\tan(h)}
\]
where $P$ is measured in $\mathrm{hPa}$, $T$ in $^\circ\mathrm{C}$, and $h$ is
the apparent altitude. Calculate the atmospheric refraction $(R)$ in
arcseconds and the true zenith distance $(z_{\mathrm{true}})$.

\medskip\noindent\textbf{Part 3 (15 points): Yersin's Lines of Position}

\begin{description}
\item[(a)] (5 points) Yersin adopted an assumed latitude of his observatory to
  be $\phi_0 = 12^\circ 15' 00''\,\mathrm{N}$, and an assumed time error
  $\Delta T_0 = 50\,\mathrm{s}$. To determine how small errors in latitude
  $(\mathrm{d}\phi)$ and clock time $(\mathrm{d}T)$ affect the observed zenith
  distance, we start with the spherical law of cosines for the
  Pole-Zenith-Star (PZS) triangle:
  \[
  \cos z = \sin\phi \sin\delta + \cos\phi \cos\delta \cos H \quad (1)
  \]
  Taking the differential of this equation with respect to latitude $\phi$ and
  hour angle $H$ (recognizing that $\mathrm{d}H = -15\cos\phi\,\mathrm{d}T$
  due to Earth's rotation), show that a change in latitude and time produces a
  shift in the calculated zenith distance given by:
  \[
  \mathrm{d}z = -\sin A \cos\phi\,\mathrm{d}H - \cos A\,\mathrm{d}\phi
  \]
  Using the definition of the intercept
  $p = z_{\mathrm{calc}} - z_{\mathrm{true}} = -\mathrm{d}z$
  ($z_{\mathrm{calc}}$ is the calculated and $z_{\mathrm{true}}$ the true
  zenith distances) and defining spatial corrections
  $x = \cos\phi\,\mathrm{d}H$ (East-West clock correction in arcseconds) and
  $y = \mathrm{d}\phi$ (North-South latitude correction in arcseconds), the
  Line of Position (LOP) equation is defined as:
  \[
  p = x\sin(A) + y\cos(A)
  \]
  Hints: Taking the derivative of equation $(1)$ gives:
  \[
  -\sin z\,\mathrm{d}z
  = (\cos\phi \sin\delta - \sin\phi \cos\delta \cos H)\,\mathrm{d}\phi
    - (\cos\phi \cos\delta \sin H)\,\mathrm{d}H
  \]
\item[(b)] (10 points) After reducing the raw data, the following table is
  obtained.

  \begin{center}
  \begin{tabular}{ccc}
  \toprule
  Star & Azimuth ($A$) & Intercept ($p$) \\
  \midrule
  A & $285.0^\circ$ & $-0.40''$ \\
  B & $245.0^\circ$ & $-18.77''$ \\
  C & $70.0^\circ$ & $+18.08''$ \\
  D & $125.0^\circ$ & $-8.15''$ \\
  \bottomrule
  \end{tabular}
  \end{center}

  Plot the four LOPs on $(x, y)$ plane (mark your graph as ``D02.3b'').
  Graphically determine the center of the common intersection region
  $(x_0, y_0)$.
\end{description}

\medskip\noindent\textbf{Part 4 (17 points): Recovering the Location of Yersin Observatory and Error Analytics}

\begin{description}
\item[(a)] (4 points) Determine the true latitude of the observatory
  $(\phi_{\mathrm{true}})$ and the true time error $(\Delta T_{\mathrm{true}})$
  using the time correction $\mathrm{d}T_{\mathrm{corr}} = x_0/(15\cos\phi_0)$
  (in seconds). Assumed that $\mathrm{d}\phi = \phi_{\mathrm{true}} - \phi_0$.
\item[(b)] (8 points) Calculate the systematic error $(\mathrm{d}z)$ of the
  astrolabe (in arcseconds) using the East-West time difference; assuming that
  the average azimuth for the East stars is $90^\circ$ and the West stars is
  $270^\circ$.
\item[(c)] (2 points) Calculate the observation random error
  $\varepsilon = |x_0\sin(A) + y_0\cos(A) - p_{\mathrm{obs}}|$ for the four
  stars.
\item[(d)] (3 points) Calculate the position error $d$ (in meters) between the
  observer's assumed coordinates and the calculated true position
  $(x_0, y_0)$. Assume $1''$ of angular distance corresponds to
  $31\,\mathrm{m}$ on the Earth's surface.
\end{description}
